% Generated draft: requires mathematical and source review.
\documentclass[11pt]{article}
\usepackage[margin=25mm]{geometry}
\usepackage{fontspec}
\setmainfont{Latin Modern Roman}
\usepackage{amsmath,amssymb,mathtools,unicode-math}
\setmathfont{Latin Modern Math}
\usepackage{enumitem,needspace,xurl,hyperref}
\setlength{\parindent}{0pt}
\setlength{\parskip}{5pt}
\setlength{\emergencystretch}{3em}
\newcommand{\N}{\mathbb{N}}
\newcommand{\Z}{\mathbb{Z}}
\newcommand{\Q}{\mathbb{Q}}
\newcommand{\R}{\mathbb{R}}
\newcommand{\C}{\mathbb{C}}
\newcommand{\F}{\mathbb{F}}
\newcommand{\A}{\mathbb{A}}
\newcommand{\PP}{\mathbb{P}}
\newcommand{\E}{\mathbb{E}}
\newcommand{\eps}{\varepsilon}
\newcommand{\dd}{\,\mathrm{d}}
\begin{document}
\section*{UnsolvedMath ID: 801150}
\textbf{Source problem code:} catalogue entry
\subsection{Definitions and mathematical statement}
The supplied record poses a research programme rather than a single conjecture with specified hypotheses. The following conventional notation explains its terminology without resolving the missing choices.

In a standard local potential presentation of the Witten--Dijkgraaf--Verlinde--Verlinde (WDVV) equations, let \(n\geq 1\), let \(\eta=(\eta_{ij})\) be a constant, invertible, symmetric matrix, and write \((\eta^{ij})=\eta^{-1}\). For a potential \(F\) in local coordinates \(t^1,\ldots,t^n\), with the required third derivatives, set
\[
F_{ijk}=\frac{\partial^3 F}{\partial t^i\partial t^j\partial t^k}.
\]
The equations are
\[
\sum_{a,b=1}^{n}F_{ija}\eta^{ab}F_{bk\ell}
=
\sum_{a,b=1}^{n}F_{ika}\eta^{ab}F_{bj\ell},
\qquad 1\leq i,j,k,\ell\leq n.
\]
Equivalently, these equations express associativity of the multiplication with structure constants
\[
c_{ij}^{\ k}=\sum_{a=1}^{n}\eta^{ka}F_{ija}.
\]
This is contextual notation: the record does not specify the ground field, regularity class, coordinate domain, metric, unit normalization, or additional conditions relating these equations to its intended Frobenius manifolds.

For an endomorphism \(A\) of an \(n\)-dimensional vector space, the usual meaning of \(\mathrm{gl}\)-regularity is
\[
\deg m_A=n,
\]
where \(m_A\) is the minimal polynomial; equivalently, the minimal and characteristic polynomials coincide. For a tangent bundle endomorphism \(L\), this condition can be imposed on \(L_p\) at each point \(p\) of a specified domain. The record does not specify that domain or whether regularity is required everywhere or only on a generic locus. It also does not define the pair \((L,e)\), its compatibility conditions, or its construction from a Frobenius manifold; these cannot be recovered from the statement alone.

\textbf{Research questions (Bolsinov, Konyaev, Matveev, 2024, Problem 3).}
\begin{enumerate}
\item Explain why the WDVV equations are interesting and identify their applications.
\item Determine whether the part of WDVV theory relevant to those applications can be reformulated by replacing a solution of the equations with a pair \((L,e)\), at least in the \(\mathrm{gl}\)-regular case.
\item Determine which pairs \((L,e)\) occur in known examples of Frobenius manifolds.
\end{enumerate}
The intended applications, the meaning of replacement or equivalence, the admissible pairs, and the family of known examples remain unspecified. Thus the supplied evidence does not determine a precise correspondence theorem or classification problem. The registry labels the problem open; that label is supplied metadata, not an independently verified mathematical assessment.
\subsection{Short English statement}
The problem asks what makes the WDVV equations useful, whether their application-relevant content can be expressed through pairs \((L,e)\), at least in the \(\mathrm{gl}\)-regular case, and which such pairs arise from known Frobenius manifolds. The supplied record leaves the pairs, correspondence, applications, and examples undefined, so a precise theorem statement requires further source context.
\subsection{Sources}
\begin{itemize}
\item Local UnsolvedMath registry: permanent ID \(801150\), exact source problem code \texttt{catalogue entry}, title \emph{Nijenhuis Geometry and Integrable Systems --- Problem 3}. The supplied record attributes the problem to Bolsinov, Konyaev, and Matveev, with proposed year \(2024\), and labels its status \texttt{open}.
\item Bolsinov, Konyaev, Matveev, \emph{Research Problems on Nijenhuis Geometry and Integrable Systems} (2024), Problem 3. Supplied source URL: \url{https://arxiv.org/abs/2410.04276}. The provenance reports access on \texttt{2026-10-16} and extraction by \texttt{arxiv\_latex}. These are supplied provenance statements; the source has not been independently consulted or verified for this draft.
\item Supplied upstream revision: \texttt{d371ee917ac9647fbb06cd6fc529deec6f3dda82}. The authoritative-source link is supplied, but the record contains no surrounding definitions or supporting passages specifying the intended correspondence.
\end{itemize}
\end{document}
